% ===================================================================
%  TABULKA č. 11 — Město a jeho budovy. — Požár.
%  Traduction 2026 d'apres texto/io, controlee sur texto/fr, texto/ru
%  et texto/uk. Consigne : voir texto/cs/10-tabulka-01.tex.
%
%  OUVERTURE DE LA TROISIEME SERIE : « TŘETÍ SÉRIE ».
%
%  LE TABLEAU DE LA VILLE EST LE PLUS INTERNATIONAL DE LA COLONNE,
%  comme il l'etait en galicien et pour la meme raison : la ville du
%  XIXe siecle a ete BATIE d'apres un plan, et le plan avait un
%  vocabulaire. « Bulvár », « prefektura », « lyceum », « muzeum »,
%  « tramvaj », « burza », « orchestr », « peristyl », « opera »,
%  « reportér », « komisař ».
%
%  MAIS LE TCHEQUE Y RESISTE PLUS QUE TOUTES LES AUTRES COLONNES, ET
%  C'EST L'HISTOIRE QUI LE VEUT. La renaissance nationale du XIXe
%  siecle a REFAIT ce vocabulaire a dessein, contre l'allemand :
%  « radnice » l'hotel de ville et non « ratuz », « nádraží » la gare,
%  « spořitelna » la caisse d'epargne, « hasiči » les pompiers,
%  « zvonice » le beffroi, « celnice » la douane, « soud » le
%  tribunal. Le mot y a donc l'age de son objet — regle du tableau 14
%  irlandais — mais AUSSI l'age de la decision politique qui l'a
%  fabrique.
% ===================================================================

%%K t11-apar-1 apar
\VUsaut{2.0mm}
\VUcentre{13.2pt}{90}{TŘETÍ SÉRIE}
\VUsaut{3.0mm}
\VUfilet{16mm}
\VUsaut{5.6mm}
\VUtitre{15.0pt}{60}{TABULKA č. 11}
\VUsaut{1.4mm}
\VUpk{11.4pt}{40}{Město a jeho budovy. --- Požár.}
\VUsaut{2.2mm}

%%K t11-tit-1 sub
\VUpk{10.2pt}{40}{Předměstí.}
\VUsaut{3.0mm}

%%K t11-c1-01-1 p
1. --- Bydlím ve velkém městě ležícím sto kilometrů od oceánu, které je přesto prvořadým
\VUgras{přístavem} \textsuperscript{(1)}. Řeka, která je lemuje, má šířku čtyři sta metrů a
tvoří velmi krásnou rejdu.

%%K t11-c1-02-1 p
2. --- Celá část města, která je na \VUgras{nábřežích} \textsuperscript{(2)}, má zcela námořní
a průmyslový vzhled a tvoří \VUgras{předměstí} \textsuperscript{(3)}, které obývají jen
námořníci a dělníci. Ti pracují v \VUgras{továrnách} \textsuperscript{(4)} a v
\VUgras{plynárně} \textsuperscript{(5)}, jejíž vysoký
\VUgras{komín} \textsuperscript{(6)} a \VUgras{plynojem} je vidět z velké dálky.

%%K t11-c1-03-1 p
3. --- Ve středověku bylo město opevněno a ještě se nacházejí některé zbytky jeho hradeb,
zvláště \VUgras{brána} \textsuperscript{(8)} starého \VUgras{opevněného hradu}
\textsuperscript{(7)}, po stranách se dvěma \VUgras{věžemi} \textsuperscript{(9)}, které dnes
slouží jako \VUgras{vězení.}

%%K t11-c1-04-1 p
4. --- V celé té \VUgras{čtvrti,} která vypadá velmi staře, je jen jedna budova poměrně
moderní: je to \VUgras{celnice} \textsuperscript{(10)}. Leží mezi nábřežím a
\VUgras{uličkou} \textsuperscript{(11)}, která vede ke staré \VUgras{radnici}
\textsuperscript{(12)}. Tuto poslední budovu lze snadno poznat podle obrovské
\VUgras{zvonice} \textsuperscript{(13)}, která ovládá staré město.

%%K t11-c1-05-1 p
5. --- Na druhé straně ulice \VUgras{stojí} budova postavená ve stejném slohu jako celnice: to
je \VUgras{burza} \textsuperscript{(14)}. Jedno z jejích průčelí hledí na malé náměstí, kde je
\VUgras{prefektura} \textsuperscript{(15)}. Naše město, ač není hlavním městem, je přece
významným hlavním městem departementu.

%%K t11-tit-2 sub
\VUsaut{5.0mm}
\VUpk{10.2pt}{40}{Nejvyšší část města.}
\VUsaut{3.4mm}

%%K t11-c1-06-1 p
6. --- Posádka, dosti početná, je ubytována v prostorných \VUgras{kasárnách}
\textsuperscript{(16)} velmi zdravých; táhnou se až k mříži \VUgras{městského parku}
\textsuperscript{(17)}. \VUgras{Bulvár} \textsuperscript{(19)}, který vede k tomu parku,
prochází za \VUgras{spořitelnou} \textsuperscript{(18)} a ústí na náměstí
\VUgras{katedrály} \textsuperscript{(20)}.

%%K t11-c1-07-1 p
7. --- Všechny ty budovy se rozkládají v podobě amfiteátru na malém pahorku, kde je také naše
prostorné \VUgras{lyceum} \textsuperscript{(21)}.

%%K t11-c1-07-2 p
Sestoupí-li se cestou poněkud svažitou, která se připojuje k Velké ulici, dojde se k
\VUgras{soudu} \textsuperscript{(22)}, před nímž se rozkládá příjemný \VUgras{veřejný
sad} \textsuperscript{(26)}. Ten \VUgras{parčík} není příliš velký, ale tvoří uprostřed města
oázu zeleně a chladu. Proto jej velmi navštěvují obyvatelé města, kteří rádi přicházejí
usednout pod plachtu \VUgras{kavárny} \textsuperscript{(25)}, aby naslouchali vojenským
hudebníkům, kteří ve čtvrtek a v neděli hrají v
\VUgras{altánu} \textsuperscript{(27)} postaveném mezi trávníky a květinovými záhony.

%%K t11-c1-08-1 p
8. --- Na jednom z těch trávníků stojí bronzová \VUgras{socha} \textsuperscript{(28)}, pomník
vztyčený na památku jednoho z našich spoluobčanů, který prokázal velké služby svému rodnému
městu.

%%K t11-tit-3 sub
\VUsaut{4.0mm}
\VUpk{10.2pt}{40}{Dopravní prostředky.}
\VUsaut{3.4mm}

%%K t11-c1-09-1 p
9. --- Naše město dosud není osvětleno elektrickým světlem, má jen \VUgras{plynové lampy}
\textsuperscript{(29)}. Ale \VUgras{zdejší noviny} vedou kampaň, aby byla ta soustava
osvětlení zlepšena, \VUgras{jako} už byla zdokonalena \VUgras{soustava pohonu} tramvají.
\VUgras{Staré} vozy, \VUgras{tažené} \VUgras{koňmi,} byly prakticky nahrazeny
\VUgras{elektrickými tramvajemi} \textsuperscript{(31)}, které rychle dopravují cestující z
jednoho konce \VUgras{města na druhý} za velmi \VUgras{levnou} cenu.

%%K t11-c1-10-1 p
10. --- Blízko soudu je \VUgras{banka} \textsuperscript{(32)}, kde se eskontují obchodní cenné
papíry. \VUgras{Bankovní poslové} \textsuperscript{(33)} chodí vybírat pohledávky do domů.

%%K t11-tit-4 sub
\VUsaut{4.0mm}
\VUpk{10.2pt}{40}{Umění a vzdělání.}
\VUsaut{3.4mm}

%%K t11-c1-11-1 p
11. --- Krásná budova, která stojí blízko, je \VUgras{muzeum.} Obsahuje zároveň městskou
\VUgras{knihovnu} \textsuperscript{(34)}, krásné \VUgras{přírodovědné sbírky}
\textsuperscript{(35)} (Přírodovědné muzeum) a půvabnou \VUgras{obrazárnu}
\textsuperscript{(36)}, která tvoří skvělou stálou \VUgras{výstavu.}

%%K t11-c1-11-2 p
Otevírá se veřejnosti dvakrát týdně, ale cizinci ji smějí navštívit kterýkoli den za spropitné
\VUgras{hlídači} \textsuperscript{(37)}. \VUgras{Malíři} \textsuperscript{(38)} tam přicházejí
pracovat. Je vidět je \VUgras{před} jejich stojanem, s \VUgras{paletou} v ruce, jak kopírují
slavné obrazy.

%%K t11-c1-12-1 p
12. --- Na výšině, za muzeem, se zahlédne \VUgras{kupole} \textsuperscript{(40)}
\VUgras{nemocnice} \textsuperscript{(39)} a budovy \VUgras{učitelského ústavu}
\textsuperscript{(41)}, v němž byli školeni učitelé našich obecných škol. Na Divadelním
náměstí je \VUgras{poštovní úřad} \textsuperscript{(43)} zřízený v \VUgras{soukromém domě}
\textsuperscript{(42)}.

%%K t11-tit-5 sub
\VUsaut{5.0mm}
\VUpk{11.4pt}{40}{Požár.}
\VUsaut{3.0mm}

%%K t11-tit-6 sub
\VUpk{10.2pt}{40}{Volání o požáru.}
\VUsaut{3.4mm}

%%K t11-c2-01-1 p
1. --- Městem se šíří strašlivá zpráva: v divadle právě vypukl požár! V okamžiku, kdy
přicházím na \VUgras{náměstí} \textsuperscript{(96)}, je zástup už shromážděn na
\VUgras{chodníku} \textsuperscript{(46)} a na \VUgras{vozovce} \textsuperscript{(47)}.
\VUgras{Poštovní zřízenci} \textsuperscript{(44)} a \VUgras{listonoši}
\textsuperscript{(45)} vycházejí z úřadu. \VUgras{Řidič tramvaje} \textsuperscript{(48)} musel
zastavit svůj vůz, aby nechal projet městský
\VUgras{sanitní vůz} \textsuperscript{(50)}, který přijíždí s \VUgras{chirurgem}
\textsuperscript{(52)} a s \VUgras{ošetřovatelem} \textsuperscript{(51)}, aby ošetřili
\VUgras{raněného} \textsuperscript{(53)} přineseného na \VUgras{nosítkách}
\textsuperscript{(54)} blízko \VUgras{novinového stánku} \textsuperscript{(55)}.

%%K t11-c2-02-1 p
2. --- Rota pěchoty, která se vracela z cvičení, se zastavila, aby zajistila pořádkovou službu
s \VUgras{městskými strážníky na koních} \textsuperscript{(61)}.
\VUgras{Jezdci,} jejichž koně se vzpínají, to dokážou lépe než \VUgras{vojáci}
\textsuperscript{(57)} nebo \VUgras{četníci} \textsuperscript{(63)}, aby zahnali zpět
\VUgras{zvědavce} \textsuperscript{(56)}. Ostatně četníci jsou velmi zaneprázdněni.
Jeden z nich ukazuje \VUgras{policistům} \textsuperscript{(64)}, kam je třeba odnést další
\VUgras{oběť} \textsuperscript{(65)}, statečného hasiče spadlého ze žebříku.

%%K t11-c2-03-1 p
3. --- \VUgras{Důstojník} \textsuperscript{(66)} přesně určuje místo, kam se má postavit
\VUgras{parní stříkačka} \textsuperscript{(67)}, která přijíždí. Zatímco čeká na ten mohutný
stroj, dal postavit do palebného postavení \VUgras{ruční stříkačku} \textsuperscript{(68)},
jejíž dlouhá \VUgras{hadice} \textsuperscript{(69)} chrlí na oheň proudy vody. Šest hasičů ji
obsluhuje, zatímco jejich druh, který drží
\VUgras{proudnici} \textsuperscript{(70)}, míří \VUgras{proud}
\textsuperscript{(71)} do středu požáru.

%%K t11-c2-04-1 p
4. --- Na druhé straně divadla opřeli velký \VUgras{žebřík} \textsuperscript{(72)}. Umožňuje
hasičům přiblížit se k plamenům, které šlehají mezi víry černého
\VUgras{dýmu}
\textsuperscript{(75)}.

%%K t11-tit-7 sub
\VUsaut{5.0mm}
\VUpk{10.2pt}{40}{Statečné činy.}
\VUsaut{3.4mm}

%%K t11-c2-05-1 p
5. --- \VUgras{Požár} \textsuperscript{(74)} vznikl ve foyeru \VUgras{divadla}
\textsuperscript{(73)}, zatímco se zkoušela \VUgras{opera}, jejíž první
\VUgras{představení} se mělo konat zítra. Naštěstí bylo v hledišti jen málo
\VUgras{diváků} \textsuperscript{(76)}. Nešťastníci se uchýlili na \VUgras{balkon}
\textsuperscript{(77)}, kam jim stateční \VUgras{hasiči} \textsuperscript{(86)} přicházejí na
pomoc s žebříkem a s lany.

%%K t11-c2-06-1 p
6. --- Pod \VUgras{peristylem} \textsuperscript{(78)}, kde \VUgras{plakát}
\textsuperscript{(79)} oznamuje představení, je vidět \VUgras{hudebníky}
\textsuperscript{(81)} \VUgras{orchestru} \textsuperscript{(80)} prchat a odnášet své
nástroje: \VUgras{housle} \textsuperscript{(81)}, \VUgras{flétnu} \textsuperscript{(82)},
\VUgras{violoncello} \textsuperscript{(83)},
\VUgras{trombon} \textsuperscript{(84)}, \VUgras{buben} \textsuperscript{(85)} atd.
Snaží se zachránit část \VUgras{nábytku} \textsuperscript{(87)}.

%%K t11-c2-07-1 p
7. --- \VUgras{Herci} \textsuperscript{(89)} a \VUgras{herečky} \textsuperscript{(88)},
\VUgras{zpěváci} a \VUgras{zpěvačky} museli prchnout ve svém divadelním \VUgras{kostýmu}
\textsuperscript{(90)}. Tenor vysvětluje
\VUgras{novináři} \textsuperscript{(92)}, jak se to stalo. \VUgras{Reportér} přiběhl
hned od prvního okamžiku s úřady: \VUgras{prefektem} \textsuperscript{(91)},
\VUgras{starostou} \textsuperscript{(93)}, \VUgras{policejním komisařem}
\textsuperscript{(94)} a \VUgras{soudním úředníkem} \textsuperscript{(95)}, kteří povedou
vyšetřování, aby soudili o odpovědnosti.
